\chapter{Drei Fragen über die christliche Kirche}
\quellseite{3}

\section[3. Mit welchen Kirchen sollen wir …]{3. Mit welchen Kirchen sollen wir Gemeinschaft halten und von welchen sollen wir uns trennen?}

\subsection[b. Eine rechtgläubige Kirche soll …]{b. Eine rechtgläubige Kirche soll sich von falschgläubigen Kirchen trennen}

Wir haben schon früher gesehen, daß eine falschgläubige Kirche eine solche Kirche ist, welche ein falsches, verfälschtes Glaubensbekenntniß hat, welche sich nicht in allen Stücken zu Gottes reinem Wort und Sacrament hält und nicht in allen Stücken den rechten und wahren christlichen Glauben bekennt.

Falschgläubig sind hiernach erstlich alle Kirchen, die jedem rechtgläubigen Christen gegenüber sich selbst frei offen als falschgläubig darstellen durch ihr, daß ich so sage, officielles Bekenntniß zu falscher Lehre.

Es ist die Pflicht jedes rechtgläubigen Christen, diese zu wissen und zu kennen. Denn der Apostel sagt durch den Heiligen Geist: „Ich ermahne aber euch, lieben Brüder, daß ihr aufsehet auf die, die da Zertrennung und Aergerniß anrichten neben der Lehre, die ihr gelernet habt, und weichet von denselbigen“, Röm. 16, 17. Und der HErr JEsus sagt: „Sehet euch vor vor den falschen Propheten“, Matth. 7, 15. Und die falschgläubigen Kirchen, die selbst das Schild des falschen Bekenntnisses ganz sichtbar aushängen, die kann man doch leicht kennen und auf sie aufsehen und sich vor ihnen vorsehen. Welche Entschuldigung könnte ein rechtgläubiger Christ haben, wenn er die nicht kennte und von denen betrogen würde?

\quellseite{4}

Solche falschgläubige Kirchen sind außer der Kirche des Antichristen, der römisch-katholischen, auch noch deren Schwester, die griechisch-katholische Kirche, und die reformirte Kirche mit den zahllosen ihr verwandten Secten, sowie die unirte Kirche, die Ja und Nein, Wahrheit und Lüge friedlich unter Einen Hut bringen will.

Falschgläubig sind aber zum andern auch alle Kirchen, welche zwar, daß ich wiederum den Ausdruck gebrauche, \gesperrt{officieller Weise, nämlich in ihrer Constitution, das rechte und unverfälschte Glaubensbekenntniß als das ihrige bezeichnen und auch den Namen der rechtgläubigen Kirche führen}; welche sich aber dennoch nicht in allen Stücken zu Gottes reinem Wort und Sacrament halten und nicht in der That und wahrhaftig in allen Stücken den rechten und wahren christlichen Glauben bekennen.

Und unter diesen möchte ich wiederum einen zweifachen Unterschied machen.

Unter diesen finden sich nämlich einmal solche, in welchen mehr oder weniger grobe Lehrwillkür herrscht, das heißt, in welchen jeder glaubt, lehrt und bekennt, was er will. — Solche Lehrwillkür herrscht auf allergröbste Weise in den „lutherischen“ Landeskirchen Deutschlands. Da weiß ich, um nur Ein Beispiel anzuführen, nicht einen einzigen theologischen Professor, der wirklich in allen Stücken die rechte Lehre führt; wohl aber weiß ich eine ganze große Anzahl, die baren Unglauben lehrt. Und die unterrichten die zukünftigen Pastoren!

Und dann finden sich unter diesen solche, die \gesperrt{ganz sonderlich in gewissen Stücken sich zu falscher Lehre bekennen und dieselbe führen und vertheidigen und die rechte Lehre verwerfen und verlästern}. — Solches geschieht zum Exempel hierzulande von den Synoden von Ohio und von Iowa in den Lehren von der Bekehrung und der Gnadenwahl.

Auch derartige falschgläubige Kirchen, die sich mit dem Namen der rechtgläubigen Kirche schmücken und das Bekenntniß derselben als Aushängeschild gebrauchen, muß ein rechtgläubiger Christ wissen und kennen, damit er sein Kleinod, die rechte Lehre, und irgend ein Kleinod der rechten Lehre, nicht verliere.

Von allen falschgläubigen Kirchen soll eine rechtgläubige Kirche sich trennen.

In welcher Weise?

Eine rechtgläubige Kirche soll sich von den falschgläubigen Kirchen erstlich und vor allen Dingen in der Weise trennen, daß sie \gesperrt{auf Lehrzucht hält}, das heißt, daß sie streng darauf sieht, daß in ihrer Mitte nur die rechte und reine Lehre des göttlichen Wortes gelehrt wird, und daß sie eben gerade die von den falschgläubigen Kirchen angefochtenen und verfälschten und verworfenen Lehren des göttlichen Wortes — daß sie eben die laut und öffentlich bekennt, aus dem Worte Gottes mit klaren Gründen beweist und vertheidigt und die entgegenstehenden falschen Lehren, wiederum aus dem Worte Gottes, widerlegt, verwirft, verdammt, item, bekämpft.

Solches muß allezeit die erste und vornehmste Weise sein, in der eine rechtgläubige Kirche von den falschgläubigen Kirchen sich trennt.

Geschieht dies, wie es denn soll, mit rechtem heiligen Ernst, so kommt es von selbst, und soll kommen, daß eine rechtgläubige Kirche sich auch äußerlich von den falschgläubigen Kirchen absondert. Denn das unter einem andern Panier Stehen und das Bekämpfen fordert Absonderung.

Und wie soll dies Absondern geschehen?

So:

\gesperrt{Eine rechtgläubige Kirche soll mit falschgläubigen Kirchen keine Kirchengemeinschaft haben.}

Was heißt das?

Das heißt zum ersten, daß eine rechtgläubige Kirche mit falschgläubigen Kirchen keine Gemeinschaft haben soll in den Werken, welche die eigentlichen Werke der Kirche sind. Und das heißt zum andern, daß eine rechtgläubige Kirche mit falschgläubigen Kirchen in gar keinen Werken, wo die Kirchen als Kirchen in Anspruch genommen werden, Gemeinschaft haben soll.

Also eine rechtgläubige Kirche soll mit falschgläubigen Kirchen keine Gemeinschaft haben in der Predigt des Wortes, in der Verwaltung der Sacramente, item, im öffentlichen Gottesdienst irgend welcher Art. Das sind die eigentlichen Werke der Kirche.

Eine rechtgläubige Kirche darf daher zum Beispiel falschgläubige Prediger nicht auf ihren Kanzeln und ihre Prediger nicht auf falschgläubigen Kanzeln predigen lassen.\footnote{Etwas anderes wäre es, wenn eine falschgläubige Kirche einen rechtgläubigen Prediger bäte, zu ihr zu kommen und ihr die Lehre seiner, der rechtgläubigen Kirche darzulegen. Das könnte er thun. Nur müßte er dann auch wirklich als Zeuge der Wahrheit auftreten und durchaus den Schein meiden, als ob es Theilnehmen am falschgläubigen Gottesdienst sei.} Sie darf ihre Glieder nicht gastweise am Sacrament der Falschgläubigen, noch Falschgläubige gastweise an ihrem Sacrament Theil nehmen lassen. Sie darf nicht gemeinsam mit Falschgläubigen Mission treiben. Sie darf nicht gemeinsam mit Falschgläubigen Gottesdienste und Gebetsversammlungen veranstalten. Sie darf nicht leiden, daß ihre Prediger gemeinsam mit falschgläubigen Predigern bei Beerdigungen amtiren. Sie darf Kirchen- oder Synodalversammlungen der Falschgläubigen nicht durch Deputirte beschicken, noch ihre Kirchenversammlungen durch Deputirte falschgläubiger Kirchen beschicken lassen. Und was desgleichen mehr ist.

Ebenso darf eine rechtgläubige Kirche mit falschgläubigen Kirchen keine Gemeinschaft haben, wenn irgend ein öffentliches Wohl gefördert oder irgend ein öffentliches Uebel abgestellt werden soll — sofern da eben die Kirchen als Kirchen in Anspruch genommen werden und nicht nur die einzelnen Glieder der Kirchen als Glieder der menschlichen Gesellschaft und als Staatsbürger in Betracht kommen.

Eine rechtgläubige Kirche darf daher zum Beispiel nicht in Gemeinschaft und Verbindung mit falschgläubigen Kirchen Maßregeln berathen und ergreifen gegen das Laster der Trunkenheit oder das der Entheiligung des Sonntags, oder für Beförderung der öffentlichen Sittlichkeit etc.

Daß eine rechtgläubige Kirche sich so von falschgläubigen Kirchen trennen und sich von ihnen absondern soll, das lehrt Gottes Wort.

Wenn der Heilige Geist sagt: „Ein wenig Sauerteig versäuert den ganzen Teig“, Gal. 5, 9., und damit vor
\quellseite{5}
der geringsten falschen Lehre und vor jedem Anfang und Ansatz derselben warnt, will er dann nicht, daß eine rechtgläubige Kirche in ihrer Mitte Lehrzucht übt? Und wenn er befiehlt, daß jeder Prediger „halte ob dem Wort, daß gewiß ist und lehren kann, auf daß er mächtig sei, zu ermahnen durch die heilsame Lehre, und zu strafen die Widersprecher“, Tit. 1, 9., will er dann nicht, daß die rechtgläubige Kirche erstens Lehrzucht übe und zum andern auch eben gerade die von den falschgläubigen Kirchen angefochtenen und verfälschten und verworfenen Lehren des göttlichen Wortes laut und öffentlich bekenne, aus dem Worte Gottes mit klaren Gründen beweise und vertheidige und die entgegenstehenden falschen Lehren widerlege, verwerfe, verdamme, bekämpfe? Wie? Meinst du, daß ich zu hart rede? Kannst du das meinen, wenn du liest, daß Paulus durch den Heiligen Geist schreibt: „Aber so auch wir, oder ein Engel vom Himmel, euch würde Evangelium predigen anders, denn das wir euch geprediget haben, der sei verflucht“? Ja, damit jeder erkenne, daß er wirklich aus dem Heiligen Geist, und nicht etwa in der Hitze und Uebereilung so geschrieben habe, fährt er fort: „Wie wir jetzt gesagt haben, so sagen wir auch abermal: So jemand euch Evangelium predigt anders, denn das ihr empfangen habt, der sei verflucht“, Gal. 1, 8. 9.

Und wenn der Apostel durch den Heiligen Geist sagt: „Ich ermahne aber euch, lieben Brüder, daß ihr aufsehet auf die, die da Zertrennung und Aergerniß anrichten neben der Lehre, die ihr gelernet habt; und weichet von denselbigen“, Röm. 16, 17.; und wenn er sagt: „Einen ketzerischen Menschen meide, wenn er einmal und abermal ermahnet ist“, Tit. 3, 10.; und wenn ein anderer Apostel durch denselben Heiligen Geist schreibt: „So jemand zu euch kommt, und bringet diese Lehre nicht, den nehmet nicht zu Hause, und grüßet ihn auch nicht; denn wer ihn grüßet, der macht sich theilhaftig seiner bösen Werke“, 2 Joh. 10. 11. — will dann der Heilige Geist nicht, daß eine rechtgläubige Kirche sich auch äußerlich, was alle Kirchengemeinschaft anlangt, von den falschgläubigen Kirchen absondere?

Ei, mein Leser, wenn du meinst, daß ich zu viel gesagt habe, so gib doch der Schrift schuld; denn ich habe nur gesagt, was sie sagt.

Und das von Gott geforderte Bekenntniß zum Worte Christi und zu allen Worten Christi fordert, daß eine rechtgläubige Kirche sich in der angezeigten Weise von den falschgläubigen Kirchen trenne. Denn wie könnte eine rechtgläubige Kirche alle Worte Christi recht bekennen, wenn sie nicht die Kirchengemeinschaft aufhübe mit denen, die dem Worte Christi widersprechen? Meinst du, daß eine rechtgläubige Kirche eine rechte Bekennerin der Wahrheit sein kann, wenn sie in kirchlicher Gemeinschaft mit denen sich zeigt, die die Wahrheit verleugnen und bekämpfen?

Und wenn eine rechtgläubige Kirche in irgend einer Weise mit falschgläubigen Kirchen Kirchengemeinschaft hat, so gibt sie schweres Aergerniß. Sie schwächt dadurch ihr sonstiges Zeugniß für die Wahrheit ab. Sie bestärkt dadurch die Falschgläubigen in ihrem Irrthum. Sie fördert dadurch den so weit verbreiteten und so schädlichen kirchlichen Indifferentismus, das heißt, die Gleichgültigkeit, in welcher so viele sprechen: „Es ist einerlei, zu welcher Kirche man gehört, wenn es nur nicht die römisch-katholische ist.“ Also das göttliche Verbot, Aergerniß zu geben, muß eine rechtgläubige Kirche bewegen, sich in angezeigter Weise von falschgläubigen Kirchen zu trennen.

Und endlich: Gott will haben, daß es eine rechtgläubige Kirche gebe und daß sie für die Wahrheit Zeugniß ablege. Gott will also haben, daß eine rechtgläubige Kirche ganz ernstlich auf ihre Selbsterhaltung bedacht sei. Aber nimmermehr wird eine rechtgläubige Kirche erhalten bleiben, wenn sie sich von den falschgläubigen Kirchen nicht getrennt hält. Durch irgend welche Kirchengemeinschaft mit falschgläubigen Kirchen wird eine rechtgläubige Kirche anfaulen, und das schnell, wie ein guter Apfel schnell anfault, wenn er unter angefaulten liegt. Also auch die in diesem Falle so heilige und gottgefällige Selbsterhaltung fordert gebieterisch, daß eine rechtgläubige Kirche sich in angezeigter Weise von falschgläubigen Kirchen trenne.

Die Hauptsache aber ist und bleibt immer das, daß Gott will und in seinem Worte gebietet, daß eine rechtgläubige Kirche sich von falschgläubigen Kirchen trenne. Da gilt es, ohne Bedenken und Dreinreden zu gehorchen. C. M. Z.
